\Article{341}
Die erste Hauptaufgabe dieser Untersuchungen besteht darin, die Anzahl der Primfunktionen jedes Grades zu bestimmen. Für einen bestimmten Modul ist nämlich die Anzahl aller verschiedenen (inkongruenten) Funktionen jeden Grades endlich. Von diesen sind einige aus Primfunktionen niedrigeren Grades zusammengesetzt, andere selbst prim; daher ist auch die Anzahl der letzteren endlich. Eine strenge Entwicklung dieses Gegenstandes ist recht heikel. Wir werden mit den einfacheren Fällen beginnen.

Setzt man den Modul $=p$, so ist die Anzahl aller verschiedenen Funktionen des Grades $n$ von der Form
\[
x^n+Ax^{n-1}+Bx^{n-2}+Cx^{n-3}+\text{usw.}
\]
gleich $p^n$. Die Anzahl der Koeffizienten $\pdot A,B\pdot C$ usw. ist nämlich $n$, und jeder kann unabhängig von den übrigen $\equiv0,1,2,3\dotstwo(p-1)(\mathrm{mod.}\,p)$ sein. Aus der Kombinatorik ergeben sich daher $p^n$ verschiedene Kombinationen, die somit die Gesamtheit aller verschiedenen Funktionen dieses Grades bestimmen.

So gibt es $p$ Funktionen ersten Grades, nämlich $x,x+1,x+2$ bis $x+p-1$, ferner $pp$ Funktionen zweiten Grades usw.

\Article{342}
Wir haben schon darauf hingewiesen, dass alle Funktionen ersten Grades als prim anzusehen sind. Beschränken wir uns also, was für unser Vorhaben genügt, auf die Funktionen, deren höchster Term den Koeffizienten $1$ hat, so gibt es $p$ Funktionen ersten Grades oder einer Dimension.

\Needspace{8\baselineskip}
Alle Funktionen zweiten Grades sind entweder aus zwei Funktionen ersten Grades zusammengesetzt oder prim. Aus der Kombinatorik ist bekannt, dass sich $p$ verschiedene Dinge bei zugelassenen Wiederholungen auf $\frac{p\pdot p+1}{1\pdot2}$ verschiedene Arten kombinieren lassen. Daher gibt es ebenso viele Funktionen, die aus zwei Primfunktionen ersten Grades zusammengesetzt sind, und somit $pp-\frac{p\pdot p+1}{1\pdot2}=\frac12(pp-p)$ Primfunktionen zweiten Grades.

\SourceLeaf{230}
Ebenso sind von allen Funktionen dritten Grades, deren Anzahl $p^3$ ist, diejenigen auszuschließen, die aus drei Primfunktionen ersten Grades zusammengesetzt sind und deren Anzahl $\frac{p\pdot p+1\pdot p+2}{1\pdot2\pdot3}$ ist. Dazu kommen diejenigen, die aus einer Primfunktion ersten und einer zweiten Grades zusammengesetzt sind und deren Anzahl $p\pdot\frac12(pp-p)$ ist. Lässt man diese weg, so bleiben $\frac13(p^3-p)$ übrig. So viele Primfunktionen dritten Grades gibt es also. Offenbar kann man auf diese Weise stets fortfahren.

\Article{343}
Um diese Rechnungen leichter auszuführen und zugleich den Weg zur Entwicklung des allgemeinen Gesetzes zu bahnen, wollen wir die Sache allgemein betrachten. Der Kürze halber bezeichnen wir mit $(1)$ die Anzahl der Primfunktionen ersten Grades, mit $(2)$ die Anzahl der Primfunktionen zweiten Grades, ferner mit $(1^2)$ die Anzahl der aus zwei Primfunktionen ersten Grades zusammengesetzten Funktionen usw. usw. Allgemein bezeichnen wir mit $(1^\alpha2^\beta3^\gamma\ldots)$ die Anzahl aller aus Primfunktionen zusammengesetzten Funktionen, und zwar aus $\alpha$ ersten, $\beta$ zweiten, $\gamma$ dritten Grades usw. Ihr Grad ist also $\alpha+2\beta+3\gamma+\text{usw.}$ Aus dem Vorhergehenden und der Kombinatorik ergibt sich dann
\[
(1^\alpha2^\beta3^\gamma4^\delta\ldots)=(1^\alpha)(2^\beta)(3^\gamma)(4^\delta)\dotstwo
\]
\[
(1^\alpha)=\frac{(1)\pdot(1)+1\pdot(1)+2\pdot(1)+3\ldots(1)+\alpha-1}{1\pdot2\pdot3\pdot4\ldots\alpha}
\]
oder allgemein
\[
(a^\alpha)=\frac{(a)\pdot(a)+1\pdot(a)+2\pdot(a)+3\ldots(a)+\alpha-1}{1\pdot2\pdot3\pdot4\ldots\alpha}
\]

\Needspace{11\baselineskip}
Sammelt man schließlich alle verschiedenen Arten, die Zahl $n$ durch Addition aus den Zahlen $1,2,3,\ldots$ zusammenzusetzen, und bezeichnet diese Arten mit $\alpha\pdot1+\beta\pdot2+\gamma\pdot3+\text{usw.}$, so ist offenbar die Summe aller dieser Ausdrücke $(1^\alpha2^\beta3^\gamma\dotstwo)$ gleich der Anzahl aller Funktionen des Grades $n$, also $=p^n$. So ist
\[
\begin{aligned}
p&=(1)\\
pp&=(1^2)+(2)\\
p^3&=(1^3)+(1\pdot2)+(3)\\
p^4&=(1^4)+(1^2\pdot2)+(1\pdot3)+(2^2)+(4)\\
&\text{usw.}
\end{aligned}
\]

Offenbar geht in den Ausdruck $p^n$ außer den Größen $(1),(2),(3)$ usw. auch
\SourceLeaf{231}
$(n)$ ein. Daraus erkennt man, wie sich alle Größen aus den vorhergehenden bestimmen lassen. So findet man
\[
\begin{array}{lll}
(1)=p&(4)=\frac14(p^4-pp)&(7)=\frac17(p^7-p)\\
(2)=\frac12(pp-p)&(5)=\frac15(p^5-p)&(8)=\frac18(p^8-p^4)\\
(3)=\frac13(p^3-p)&(6)=\frac16(p^6-p^3-pp+p)&\text{usw.}
\end{array}
\]

\Article{344--346}
Aus diesem Anfang der Folge erkennt man, dass der höchste Term des Ausdrucks $(n)$ gleich $\frac1n p^n$ ist. Ist $n$ eine Primzahl, so kommt $-\frac1n p$ hinzu; ist $n$ dagegen zusammengesetzt, so tritt das Gesetz weniger deutlich hervor. Betrachtet man die Sache aber genauer, so sieht man, dass
\[
\begin{array}{ll}
p=(1)&p^5=5(5)+(1)\\
pp=2(2)+(1)&p^6=6(6)+3(3)+2(2)+(1)\\
p^3=3(3)+(1)&p^7=7(7)+(1)\\
p^4=4(4)+2(2)+(1)&p^8=8(8)+4(4)+2(2)+(1)\text{ usw.}
\end{array}
\]
wobei das Bildungsgesetz offensichtlich ist. Hat nämlich die Zahl $n$ sämtliche Teiler $\alpha,\beta,\gamma,\delta$ usw., so gilt
\[
p^n=\alpha(\alpha)+\beta(\beta)+\gamma(\gamma)+\delta(\delta)+\text{usw.}
\]

Wir gehen nun daran, die Allgemeingültigkeit dieser Beobachtung zu beweisen.

\Needspace{16\baselineskip}
\begingroup\rightskip=0pt plus1fil\relax Wir haben gezeigt, dass die Summe aller solchen Ausdrücke $(1^\alpha)(2^\beta)(3^\gamma)\ldots$, wenn stets $\alpha+2\beta+3\gamma+\ldots=n$ gilt, alle Funktionen des Grades $n$ erschöpft und somit $=p^n$ ist. Hieraus folgt, \textemdash\ \textemdash\ \textemdash. Entwickelt man
\[
\left(\frac1{1-x}\right)^{(1)}\left(\frac1{1-xx}\right)^{(2)}\left(\frac1{1-x^3}\right)^{(3)}\ldots
\]
in die Reihe $1+Ax+Bx^2\ldots=P$, so ist
\[
A=p,\quad B=p^2,\quad C=p^3\text{ usw.}
\]
\[
\frac{x\,\mathrm{d}P}{P\,\mathrm{d}x}=\frac{(1)x}{1-x}+\frac{2(2)x^3}{1-x^2}+\frac{3(3)x^3}{1-x^3}\ldots
\]
\noindent\hspace*{.28\linewidth}\textemdash\ \textemdash\ \textemdash\ \textemdash\par\endgroup

[Hieraus ergibt sich, wenn man $\frac{px}{1-px}$ für $\frac{x\,\mathrm{d}P}{P\,\mathrm{d}x}$ einsetzt und die einzelnen Brüche in unendliche Reihen entwickelt, die Richtigkeit des Satzes von selbst.]

\SourceLeaf{232}
\Article{347}
Dieser Satz lässt sich auch anders ausdrücken. Sind sämtliche Teiler der Zahl $n$ gleich $n,1,\delta,\delta',\delta'',\delta'''$ usw., so besagt der Satz, dass
\[
p^n=n(n)+(1)+\delta(\delta)+\delta'(\delta')+\text{usw.}
\]
Offenbar hat nun das Produkt aus den $(n)$ Primfunktionen des Grades $n$ den Grad $n(n)$. Entsprechendes gilt für die übrigen. Daher gilt

\emph{Das Produkt aus allen Primfunktionen ersten Grades und der Grade $n,\delta,\delta'$ usw. hat den Grad $p^n$.}

Aus diesem Satz lässt sich nun leicht der Wert des Ausdrucks $(n)$ selbst ableiten. Der Kürze halber unterdrücken wir jedoch die Rechnung, die nicht schwierig ist. Es sei also $n=a^\alpha b^\beta c^\gamma$ usw., wobei $a,b,c$ usw. verschiedene Primzahlen seien. Dann gilt
\[
n(n)=p^n-\Sigma p^{\frac na}+\Sigma p^{\frac n{ab}}-\Sigma p^{\frac n{abc}}\text{ usw.}
\]
Hier bezeichnet $\Sigma p^{\frac n{abc\dotstwo}}$ die Gesamtheit aller zu $p^{\frac n{abc\dotstwo}}$ ähnlichen Ausdrücke, wenn man die Größen $a,b,c\dotstwo$ auf jede Weise untereinander vertauscht. So ist für $n=36$ dann $36(36)=p^{36}-p^{18}-p^{12}+p^6$.

\Needspace{15\baselineskip}
Eine weitere Bemerkung sei noch erlaubt. Hat $n$ die Form $a^\alpha$ und ist $a$ prim, so ist $n(n)=p^n-p^{\frac na}$. Da $(n)$ notwendig eine ganze Zahl ist, gilt daher für jedes $p$
\[
p^n\equiv p^{\frac na}\;(\mathrm{mod.}\,n)
\]
Ist also $p$ zu $a$ teilerfremd, so gilt
\[
p^{n-\frac na}\equiv1\;(\mathrm{mod.}\,n)
\]
und für $\alpha=1$
\[
p^{a-1}\equiv1\;(\mathrm{mod.}\,a)
\]
Es ist bemerkenswert, dass sich diese Sätze auf so verschiedene Weise gewinnen lassen.
